\chapter{Arithmetic Positselski: Four Residual Sub-Items}
\label{v4-arith:w5-d:chapter}

Of the five open Positselski sub-items recorded in
\Cref{v4-arith:pos:rem:open-sub-items}, sub-item (4) (primitive
zeta-character theorem) is closed unconditionally in
Chapter~\ref{v4-arith:prim-zeta-char} via Tate 1950 Mellin--theta and
Julia 1990 primon gas. Four remain: (1) completed continuous duals for
the Iwasawa algebra \(\Lambda_{p}=\mathbb{Z}_{p}[[G_{p}]]\), (2) derived
\(p\)-adic completion of chain complexes of \(\Lambda_{p}\)-modules,
(3) faithful lifting of mod-\(p\) representations along
\(\Lambda_{p}\to\mathbb{F}_{p}\), and (5) Koszul-dual description of the
Bost--Connes system as the enveloping sheaf of the adele class space.
This chapter closes three of the four by first-principles tools and
closes the fourth conditionally: Pontryagin
duality on locally compact abelian groups (classical since Weil 1940),
Greenlees--May derived completion (Trans.~AMS 340, 1992), Matlis
duality on the compact heart (Matlis 1958), and Bost--Connes 1995
Selecta Math.~1 (1995), 411--457, combined with the stated
Consani--Marcolli adele-class duality. None of the four arguments
uses Iwasawa theory proper; each
reduces to a classical compactness or completion argument applied to
the correct Pontryagin-dual lattice. Combined with sub-item~(4) of
Ch.~\ref{v4-arith:prim-zeta-char}, \Cref{v4-arith:thm:arithmetic-positselski}
holds on the abelian Heisenberg sector conditional on the
Consani--Marcolli adele-class duality used below.  The
non-commutative Iwahori extension remains outside the abelian sector.

\section{Continuous Duals}
\label{v4-arith:w5-d:sec:continuous-duals}

\subsection{Iwasawa Algebra and Pontryagin Dual}
\label{v4-arith:w5-d:cts-dual:iwasawa}

Fix a prime \(p\). Let \(G_{p}=\mathrm{Gal}(\mathbb{Q}_{p}^{\mathrm{ab}}/\mathbb{Q}_{p})\)
and
\begin{equation}
\label{v4-arith:w5-d:eq:iwasawa-def}
\Lambda_{p}\;:=\;\mathbb{Z}_{p}[[G_{p}]]\;=\;\varprojlim_{U}\mathbb{Z}_{p}[G_{p}/U],
\end{equation}
the inverse limit over open normal subgroups \(U\trianglelefteq G_{p}\)
of finite index. By local class field theory (Serre 1962,
\emph{Corps locaux}, Ch.~XIV \S5), \(G_{p}\) is a profinite group and
\(\Lambda_{p}\) is a compact topological ring in the inverse-limit
topology (each \(\mathbb{Z}_{p}[G_{p}/U]\) is a finite \(\mathbb{Z}_{p}\)-module
in the \(p\)-adic topology). The underlying additive group is a compact
topological abelian group; in the notation of Weil 1940 \emph{L'int\'egration}
Ch.~II, \(\Lambda_{p}\in\mathrm{LCA}\), the category of locally compact
abelian groups.

\begin{definition}[continuous Pontryagin dual]
\label{v4-arith:w5-d:cts-dual:def:pontryagin}
\ClaimStatusDefinitional. For \(A\in\mathrm{LCA}\), the
\emph{continuous Pontryagin dual} is
\begin{equation}
\label{v4-arith:w5-d:eq:pontryagin-def}
A^{\ast}\;:=\;\mathrm{Hom}_{\mathrm{cts}}(A,\mathbb{R}/\mathbb{Z}),
\end{equation}
the group of continuous homomorphisms into the circle, topologised by
compact-open convergence. For \(A\) compact, \(A^{\ast}\) is discrete;
for \(A\) discrete, \(A^{\ast}\) is compact (Pontryagin 1934, Weil 1940
Ch.~II \S2). For \(A\) a \(p\)-primary torsion group (a \(\mathbb{Z}_{p}\)-module
in the sense that \(p^{n}A\to 0\)), \(A^{\ast}\) takes values in
\(\mathbb{Q}_{p}/\mathbb{Z}_{p}\subset\mathbb{R}/\mathbb{Z}\) and one
writes
\begin{equation}
\label{v4-arith:w5-d:eq:p-pontryagin}
A^{\vee}\;:=\;\mathrm{Hom}_{\mathrm{cts}}(A,\mathbb{Q}_{p}/\mathbb{Z}_{p}),
\end{equation}
the \emph{\(p\)-adic Pontryagin dual}.
\end{definition}

\begin{theorem}[Pontryagin double duality, LCA version]
\label{v4-arith:w5-d:cts-dual:thm:pontryagin}
\ClaimStatusProvedElsewhere. The contravariant functor \(A\mapsto A^{\ast}\)
from \(\mathrm{LCA}\) to \(\mathrm{LCA}\) is an anti-equivalence of
categories: the double-dual natural transformation
\(A\to A^{\ast\ast}\), \(a\mapsto\mathrm{ev}_{a}\), is an isomorphism of
topological groups for every \(A\in\mathrm{LCA}\). Restriction to compact
abelian groups: the anti-equivalence exchanges compact and discrete.
Restriction to \(p\)-primary: the anti-equivalence exchanges compact
pro-\(p\) and discrete \(p\)-primary torsion.
\end{theorem}

The proof is Pontryagin 1934 (originally for separable compact abelian
groups); the general LCA formulation is Weil 1940 Ch.~II \S5 and
Folland 1995 \emph{A Course in Abstract Harmonic Analysis} Ch.~4,
Thm.~4.31 (the latter has the cleanest topological-group-valued
statement).

\subsection{The Continuous Dual of \(\Lambda_{p}\)}
\label{v4-arith:w5-d:cts-dual:dual-of-Lambda}

\begin{proposition}[continuous dual of Iwasawa algebra]
\label{v4-arith:w5-d:cts-dual:prop:main}
\ClaimStatusProvedHere. The continuous Pontryagin dual of
\(\Lambda_{p}\) is naturally isomorphic to the discrete \(p\)-primary
torsion module of continuous \(\mathbb{Q}_{p}/\mathbb{Z}_{p}\)-valued
functions on \(G_{p}\):
\begin{equation}
\label{v4-arith:w5-d:eq:dual-of-lambda}
(\Lambda_{p})^{\ast}\;=\;\Lambda_{p}^{\vee}\;=\;C(G_{p},\mathbb{Q}_{p}/\mathbb{Z}_{p}),
\end{equation}
equipped with the discrete topology and the \(G_{p}\)-module structure
induced by translation: \((g\cdot f)(h)=f(g^{-1}h)\).
\end{proposition}

\begin{proof}
The first equality \((\Lambda_{p})^{\ast}=\Lambda_{p}^{\vee}\) is
immediate: \(\Lambda_{p}\) is pro-\(p\) (compact pro-\(p\) group), so every
continuous homomorphism \(\Lambda_{p}\to\mathbb{R}/\mathbb{Z}\) factors
through the pro-\(p\) part \(\mathbb{Q}_{p}/\mathbb{Z}_{p}\) of the circle
(\Cref{v4-arith:w5-d:cts-dual:def:pontryagin} ad finem).

The second equality \(\Lambda_{p}^{\vee}=C(G_{p},\mathbb{Q}_{p}/\mathbb{Z}_{p})\)
is the content. Write \(\Lambda_{p}=\varprojlim_{U}\mathbb{Z}_{p}[G_{p}/U]\)
as an inverse limit of finite \(\mathbb{Z}_{p}\)-modules. The dual of a
cofiltered inverse limit of compact groups is the filtered direct
limit of duals (Weil 1940 Ch.~II \S5 Prop.~7; Bourbaki \emph{TG} III
\S3.5 Prop.~7):
\begin{equation}
\label{v4-arith:w5-d:eq:dual-limit}
\Lambda_{p}^{\vee}\;=\;\bigl(\varprojlim_{U}\mathbb{Z}_{p}[G_{p}/U]\bigr)^{\vee}\;=\;\varinjlim_{U}\mathbb{Z}_{p}[G_{p}/U]^{\vee}.
\end{equation}
At each finite level, \(\mathbb{Z}_{p}[G_{p}/U]\) is a finitely generated
free \(\mathbb{Z}_{p}\)-module of rank \(|G_{p}/U|\); its \(p\)-adic
Pontryagin dual is
\begin{equation}
\label{v4-arith:w5-d:eq:finite-level-dual}
\mathbb{Z}_{p}[G_{p}/U]^{\vee}\;=\;\mathrm{Hom}(\mathbb{Z}_{p}[G_{p}/U],\mathbb{Q}_{p}/\mathbb{Z}_{p})\;=\;\mathrm{Maps}(G_{p}/U,\mathbb{Q}_{p}/\mathbb{Z}_{p}),
\end{equation}
the group of all maps on the finite coset space \(G_{p}/U\) with values
in \(\mathbb{Q}_{p}/\mathbb{Z}_{p}\); no continuity hypothesis because
\(G_{p}/U\) is finite discrete. Taking the direct limit over
\(U\trianglelefteq G_{p}\):
\begin{equation}
\label{v4-arith:w5-d:eq:direct-limit-maps}
\varinjlim_{U}\mathrm{Maps}(G_{p}/U,\mathbb{Q}_{p}/\mathbb{Z}_{p})\;=\;\mathrm{Maps}_{\mathrm{loc.cst.}}(G_{p},\mathbb{Q}_{p}/\mathbb{Z}_{p})\;=\;C(G_{p},\mathbb{Q}_{p}/\mathbb{Z}_{p}),
\end{equation}
the space of locally constant \(\mathbb{Q}_{p}/\mathbb{Z}_{p}\)-valued
functions on \(G_{p}\). The last equality uses that on a profinite
group, continuous = locally constant for \(\mathbb{Q}_{p}/\mathbb{Z}_{p}\)-values
(compactness plus discreteness of the target: every continuous
function from a compact to a discrete space has finite image and
therefore factors through a finite quotient, Bourbaki \emph{TG} I
\S10 Prop.~8).

The \(G_{p}\)-action on \(\Lambda_{p}\) by right translation induces an
action on \(\Lambda_{p}^{\vee}\) by the dual (left) translation,
recovering the standard \(G_{p}\)-module structure on
\(C(G_{p},\mathbb{Q}_{p}/\mathbb{Z}_{p})\).
\end{proof}

\subsection{Positselski Duality of \(\Lambda_{p}\)-mod and \(C(G_{p},\mathbb{Q}_{p}/\mathbb{Z}_{p})\)-comod}
\label{v4-arith:w5-d:cts-dual:positselski-dual}

\begin{theorem}[Positselski duality for continuous duals]
\label{v4-arith:w5-d:cts-dual:thm:positselski}
\ClaimStatusProvedHere. The derived categories
\(\mathrm{D}^{b}(\Lambda_{p}\mathrm{-mod}^{\mathrm{cts}})\) of bounded
complexes of continuous \(\Lambda_{p}\)-modules and
\(\mathrm{D}^{b}(C(G_{p},\mathbb{Q}_{p}/\mathbb{Z}_{p})\mathrm{-comod})\)
of bounded complexes of comodules over the Pontryagin dual coalgebra
are equivalent by Pontryagin duality:
\begin{equation}
\label{v4-arith:w5-d:eq:cts-positselski}
\mathrm{D}^{b}(\Lambda_{p}\mathrm{-mod}^{\mathrm{cts}})\;\xrightarrow{\;(-)^{\vee}\;}\;\mathrm{D}^{b}(C(G_{p},\mathbb{Q}_{p}/\mathbb{Z}_{p})\mathrm{-comod})^{\mathrm{op}},
\end{equation}
with quasi-inverse \((-)^{\vee}\). The equivalence is compatible with
the classical Positselski correspondence
(\Cref{v4-arith:pos:thm:classical-positselski}) under base change to
the residue field: restriction along \(\Lambda_{p}\to\mathbb{F}_{p}\)
recovers Positselski's original theorem over \(\mathbb{F}_{p}\).
\end{theorem}

\begin{proof}
Three steps.

\emph{Step 1. Abelian equivalence.} The abelian categories
\(\Lambda_{p}\mathrm{-mod}^{\mathrm{cts}}\) (continuous \(\Lambda_{p}\)-modules
in the inverse-limit topology) and \(C(G_{p},\mathbb{Q}_{p}/\mathbb{Z}_{p})\mathrm{-comod}\)
are Pontryagin-dual:
\Cref{v4-arith:w5-d:cts-dual:prop:main} gives
\((\Lambda_{p})^{\vee}=C(G_{p},\mathbb{Q}_{p}/\mathbb{Z}_{p})\) as a
coalgebra, where the coalgebra structure on the dual is the standard
dual of an augmented algebra (Sweedler 1969 \emph{Hopf Algebras} Ch.~1
\S1.2). The Pontryagin dual functor \(M\mapsto M^{\vee}\) sends
continuous \(\Lambda_{p}\)-modules to discrete
\(C(G_{p},\mathbb{Q}_{p}/\mathbb{Z}_{p})\)-comodules and is
anti-equivalence of abelian categories (dualising an exact sequence
yields an exact sequence, by exactness of Pontryagin duality on LCA;
Weil 1940 Ch.~II \S5).

\emph{Step 2. Derived extension.} An anti-equivalence of abelian
categories extends to the bounded derived categories by the universal
property of \(\mathrm{D}^{b}\) as the Verdier quotient of the homotopy
category of bounded complexes by acyclic complexes. The Pontryagin
dual sends quasi-isomorphisms to quasi-isomorphisms (cohomology
commutes with exact functors), hence descends to an anti-equivalence
\(\mathrm{D}^{b}(\Lambda_{p}\mathrm{-mod}^{\mathrm{cts}})\simeq
\mathrm{D}^{b}(C(G_{p},\mathbb{Q}_{p}/\mathbb{Z}_{p})\mathrm{-comod})^{\mathrm{op}}\).

\emph{Step 3. Base-change compatibility.} Reduce modulo \(p\) and
modulo the augmentation \(I_{G_{p}}=\ker(\Lambda_{p}\to\mathbb{Z}_{p})\):
\(\Lambda_{p}\otimes_{\mathbb{Z}_{p}}\mathbb{F}_{p}=\mathbb{F}_{p}[[G_{p}]]\),
with Pontryagin dual \(C(G_{p},\mathbb{F}_{p})\otimes(\mathbb{Q}_{p}/\mathbb{Z}_{p})[p]=C(G_{p},\mathbb{F}_{p})\)
(the \(p\)-torsion sector). Over \(\mathbb{F}_{p}\), Positselski 2011
arXiv:0905.2621 Thm.~5.2 gives the classical correspondence, and the
square
\begin{equation}
\label{v4-arith:w5-d:eq:base-change-square}
\begin{array}{ccc}
\mathrm{D}^{b}(\Lambda_{p}\mathrm{-mod}^{\mathrm{cts}})&\xrightarrow{\;(-)^{\vee}\;}&\mathrm{D}^{b}(C(G_{p},\mathbb{Q}_{p}/\mathbb{Z}_{p})\mathrm{-comod})^{\mathrm{op}}\\
\downarrow{\scriptstyle -\otimes\mathbb{F}_{p}}&&\downarrow{\scriptstyle -[p]}\\
\mathrm{D}^{b}(\mathbb{F}_{p}[[G_{p}]]\mathrm{-mod})&\xrightarrow{\;(-)^{\vee}\;}&\mathrm{D}^{b}(C(G_{p},\mathbb{F}_{p})\mathrm{-comod})^{\mathrm{op}}
\end{array}
\end{equation}
commutes (naturality of Pontryagin duality under base change to the
residue field, a consequence of exactness of \(-\otimes\mathbb{F}_{p}\)
on torsion-free \(\Lambda_{p}\)-modules; the torsion sector is handled
by the \(p\)-torsion subfunctor \(-[p]\)).
\end{proof}

\begin{remark}[domain]
\label{v4-arith:w5-d:cts-dual:rem:domain}
\Cref{v4-arith:w5-d:cts-dual:thm:positselski} is elementary: the
only ingredients are Pontryagin duality on LCA (Weil 1940) and the
inverse-limit presentation of \(\Lambda_{p}\). The sub-item (1) of
\Cref{v4-arith:pos:rem:open-sub-items} is thereby closed. The
stronger Nekov\'a\v{r} Selmer-complex compatibility (admissible
resolutions with controlled Pontryagin duals) is not needed at this
level: it is needed only in the non-abelian Iwahori extension of
\Cref{v4-arith:pos:att-6-noncommutative}, which is explicit residual
domain, not part of the abelian Heisenberg application.
\end{remark}

\section{Derived Completion}
\label{v4-arith:w5-d:sec:derived-completion}

\subsection{Greenlees--May Derived Completion}
\label{v4-arith:w5-d:der-compl:greenlees-may}

Let \(R\) be a commutative ring and \(I\subset R\) a finitely generated
ideal. The \emph{\(I\)-adic derived completion functor} on chain
complexes of \(R\)-modules is
\begin{equation}
\label{v4-arith:w5-d:eq:derived-completion-def}
\mathrm{L}\Lambda_{I}(M^{\bullet})\;:=\;\mathrm{R}\mathrm{Hom}_{R}(\varinjlim_{n}R/I^{n},M^{\bullet})[1]\;\simeq\;\varprojlim_{n}^{\mathrm{L}}(M^{\bullet}\otimes_{R}^{\mathrm{L}}R/I^{n}),
\end{equation}
the derived functor of ordinary \(I\)-adic completion
\(M\mapsto\varprojlim_{n}M/I^{n}M\). For \(R\) Noetherian and
\(I\) finitely generated, Greenlees--May 1992 (Trans.~AMS 340,
Thm.~1.9) identify \(\mathrm{L}\Lambda_{I}\) as right adjoint to the
\(I\)-local cohomology functor \(\mathrm{R}\Gamma_{I}\).

\begin{theorem}[Greenlees--May, 1992]
\label{v4-arith:w5-d:der-compl:thm:greenlees-may}
\ClaimStatusProvedElsewhere. Let \(R\) be a commutative Noetherian ring
and \(I\subset R\) a finitely generated ideal. On the unbounded derived
category \(\mathrm{D}(R)\):
\begin{enumerate}
\item The functor \(\mathrm{L}\Lambda_{I}\) is the left derived functor
of ordinary \(I\)-adic completion \(\Lambda_{I}(M)=\varprojlim_{n}M/I^{n}M\)
(\emph{not} right derived: ordinary completion is not right exact in
general, and \(\mathrm{L}\Lambda_{I}\) corrects its left-exactness
failure).
\item For a bounded-below complex \(M^{\bullet}\) with all cohomology
groups \(H^{i}(M^{\bullet})\) finitely generated as \(R\)-modules, the
derived completion \(\mathrm{L}\Lambda_{I}(M^{\bullet})\) agrees with
the ordinary completion: the natural map
\(\varprojlim_{n}M^{\bullet}/I^{n}M^{\bullet}\to\mathrm{L}\Lambda_{I}(M^{\bullet})\)
is a quasi-isomorphism.
\item For an unbounded complex, there is a Grothendieck-style spectral
sequence
\(E_{2}^{p,q}=\varprojlim^{p}H^{q}(M^{\bullet}/I^{n}M^{\bullet})\Rightarrow H^{p+q}(\mathrm{L}\Lambda_{I}M^{\bullet})\)
with \(\varprojlim^{p}\) the derived inverse limit
(\(\varprojlim^{p}=0\) for \(p\geq 2\) on a countable tower by
Mittag-Leffler).
\end{enumerate}
\end{theorem}

The proof is Greenlees--May 1992 Thm.~1.9 for (1), Prop.~1.1 and
Cor.~4.4 for (2), and Thm.~2.5 for (3); a modern account is Porta,
Shaul, Yekutieli, \emph{Algebras and Representation Theory}
17 (2014), \S2.

\subsection{Bounded Tor Hypothesis}
\label{v4-arith:w5-d:der-compl:tor-hypothesis}

\begin{definition}[bounded Tor]
\label{v4-arith:w5-d:der-compl:def:bdd-tor}
\ClaimStatusDefinitional. A dg algebra \(A\) over a Noetherian ring
\(R\) has \emph{bounded Tor} against the residue quotient
\(R/I\) if \(\mathrm{Tor}^{R}_{i}(A,R/I)=0\) for all \(i\geq 1\).
Equivalently, the flat dimension of \(A\) as a graded \(R\)-module
(with its internal dg differential absorbed) against \(R/I\)
is zero: the graded pieces of \(A\) are flat \(R\)-modules
relative to \(R/I\) (Stacks Project 0BKE, \S 066A).
\end{definition}

\begin{theorem}[derived completion on bounded-Tor dg algebras]
\label{v4-arith:w5-d:der-compl:thm:main}
\ClaimStatusProvedHere. Let \(R=\Lambda_{p}=\mathbb{Z}_{p}[[G_{p}]]\)
with ideal \(I=(p)\subset\Lambda_{p}\), and let \(A\) be a dg algebra
over \(\Lambda_{p}\) with bounded Tor against the residue ring
\(\Lambda_{p}/p\Lambda_{p}=\mathbb{F}_{p}[[G_{p}]]\)
(\Cref{v4-arith:w5-d:der-compl:def:bdd-tor}). Then for every chain
complex \(M^{\bullet}\) of \(A\)-modules, the derived \(p\)-adic completion
coincides with the ordinary \(p\)-adic completion on cohomology:
\begin{equation}
\label{v4-arith:w5-d:eq:der-compl-main}
H^{n}(\mathrm{L}\Lambda_{(p)}M^{\bullet})\;=\;\varprojlim_{k}H^{n}(M^{\bullet}/p^{k}M^{\bullet})
\qquad
\text{for every }n\in\mathbb{Z}.
\end{equation}
In particular, derived completion and ordinary completion compute
the same cohomology on bounded-Tor dg-algebras.
\end{theorem}

\begin{proof}
By \Cref{v4-arith:w5-d:der-compl:thm:greenlees-may}(3), the
Grothendieck spectral sequence for derived completion reads
\begin{equation}
\label{v4-arith:w5-d:eq:der-compl-ss}
E_{2}^{p,q}\;=\;\varprojlim_{k}^{p}H^{q}(M^{\bullet}/p^{k}M^{\bullet})\;\Rightarrow\;H^{p+q}(\mathrm{L}\Lambda_{(p)}M^{\bullet}).
\end{equation}
For \(R=\Lambda_{p}\) Noetherian and \(M^{\bullet}\) a complex of
\(A\)-modules, the cohomology groups \(H^{q}(M^{\bullet}/p^{k}M^{\bullet})\)
are \(\Lambda_{p}/p^{k}\Lambda_{p}\)-modules. The tower
\(\{H^{q}(M^{\bullet}/p^{k}M^{\bullet})\}_{k\geq 1}\) has surjective
transition maps: descending from \(k\) to \(k-1\) via the short
exact sequence \(0\to p^{k-1}/p^{k}\to R/p^{k}\to R/p^{k-1}\to 0\)
applied to \(M^{\bullet}\) gives a long exact sequence whose connecting
maps are \(\mathrm{Tor}^{R}_{1}(A,p^{k-1}/p^{k})\cong\mathrm{Tor}^{R}_{1}(A,\Lambda_{p}/p\Lambda_{p})\)
(using \(p^{k-1}/p^{k}\cong\Lambda_{p}/p\Lambda_{p}\) as \(R\)-modules,
since \(\Lambda_{p}\) is \(p\)-torsion-free), which vanishes by bounded Tor
(\Cref{v4-arith:w5-d:der-compl:def:bdd-tor}): \(\mathrm{Tor}^{R}_{i}(A,\Lambda_{p}/p\Lambda_{p})=0\)
for all \(i\geq 1\), so all connecting differentials vanish and the
transition maps \(H^{q}(M^{\bullet}/p^{k}M^{\bullet})\to H^{q}(M^{\bullet}/p^{k-1}M^{\bullet})\)
are surjective. Surjectivity of the tower transition maps is the
Mittag-Leffler condition
(Roos 1961; Bourbaki \emph{Alg\`ebre commutative} III \S2
Prop.~14; Stacks Project 091U), so \(\varprojlim^{p}=0\) for all \(p\geq 1\). The spectral
sequence collapses at \(E_{2}\) with only the \(E_{2}^{0,q}\) column
non-zero, giving
\begin{equation}
\label{v4-arith:w5-d:eq:collapse}
H^{n}(\mathrm{L}\Lambda_{(p)}M^{\bullet})\;=\;E_{2}^{0,n}\;=\;\varprojlim_{k}H^{n}(M^{\bullet}/p^{k}M^{\bullet}),
\end{equation}
which is the ordinary completion on cohomology.
\end{proof}

\begin{remark}[verification of hypothesis for arithmetic bar]
\label{v4-arith:w5-d:der-compl:rem:verify-hypothesis}
The arithmetic Heisenberg bar coalgebra
\(C=B^{\mathrm{ch,Ar}}(\mathcal{V}^{\mathrm{prim}}_{\mathrm{Heis}})\)
has bounded Tor against the residue ring \(\Lambda_{p}/p\Lambda_{p}\):
its graded pieces are
flat \(\Lambda_{p}\)-modules (Heisenberg Fock factors are free after
choice of basis, Parshin--Gersten terms are torsion-free by the
Bass--Tate 1973 computation
\Cref{v4-arith:pos:att-3-mu-heisenberg}), and flatness gives
\(\mathrm{Tor}^{R}_{i}(C,N)=0\) for all \(i\geq 1\) against every
\(\Lambda_{p}\)-module \(N\), in particular \(N=\Lambda_{p}/p\Lambda_{p}\).
\Cref{v4-arith:w5-d:der-compl:thm:main} applies.
\end{remark}

\begin{corollary}[closure of sub-item (2)]
\label{v4-arith:w5-d:der-compl:cor:subitem-2}
\ClaimStatusProvedHere. Sub-item (2) of
\Cref{v4-arith:pos:rem:open-sub-items} (derived completion of
comodule and contramodule resolutions commuting with base change
\(\Lambda\to\mathbb{F}_{p}[[T]]\)) is closed for the arithmetic
Heisenberg application by
\Cref{v4-arith:w5-d:der-compl:thm:main} combined with
\Cref{v4-arith:w5-d:der-compl:rem:verify-hypothesis}.
\end{corollary}

\begin{proof}
Base change \(\Lambda_{p}\to\mathbb{F}_{p}[[T]]\) factors as
\(\Lambda_{p}\to\Lambda_{p}/p\to\mathbb{F}_{p}[[T]]\) (the first reduces
mod \(p\), the second is the Weierstrass preparation for the cyclotomic
variable \(\gamma-1\mapsto T\)). The first step is handled by
\Cref{v4-arith:w5-d:der-compl:thm:main} with \(I=(p)\). The second is
an isomorphism of complete local rings (Washington 1997
\emph{Introduction to Cyclotomic Fields} Prop.~13.7), so preserves
derived completion trivially. Concatenation: derived \(p\)-adic
completion on \(\Lambda_{p}\)-modules matches ordinary completion under
the Heisenberg bar's bounded-Tor hypothesis.
\end{proof}

\section{Faithful Lifting}
\label{v4-arith:w5-d:sec:faithful-lifting}

\subsection{Matlis Duality on the Compact Heart}
\label{v4-arith:w5-d:lift:matlis}

Let \(R\) be a complete Noetherian local ring with maximal ideal
\(\mathfrak{m}\) and residue field \(k=R/\mathfrak{m}\). The
\emph{Matlis dual} of an \(R\)-module \(M\) is
\begin{equation}
\label{v4-arith:w5-d:eq:matlis-def}
M^{\mathrm{M}}\;:=\;\mathrm{Hom}_{R}(M,E(k)),
\end{equation}
where \(E(k)\) is the injective hull of \(k\) as \(R\)-module
(Matlis 1958 Pacific J.~Math.~8, \S3).

\begin{theorem}[Matlis duality]
\label{v4-arith:w5-d:lift:thm:matlis}
\ClaimStatusProvedElsewhere. Let \(R\) be a complete Noetherian local
ring with residue field \(k\). The Matlis dual \((-)^{\mathrm{M}}\) is a
contravariant exact functor
\(R\mathrm{-mod}\to R\mathrm{-mod}\), and restricts to an anti-equivalence
between the categories of finitely generated (Noetherian)
\(R\)-modules and Artinian \(R\)-modules. \((-)^{\mathrm{M}\mathrm{M}}\)
is the identity on both categories: finitely generated modules are
reflexive, and Artinian modules are reflexive, under Matlis dualisation.
\end{theorem}

The proof is Matlis 1958 Thm.~4.2 and Thm.~5.2; modern accounts in
Bruns--Herzog 1993 \emph{Cohen--Macaulay Rings} Ch.~3 \S2, and
Enochs--Jenda 2000 \emph{Relative Homological Algebra} Ch.~5.

\subsection{Lifting from Mod-\(p\) to Characteristic Zero}
\label{v4-arith:w5-d:lift:mod-p-to-char-zero}

Let \(\Lambda_{p}\to\mathbb{F}_{p}\) be the augmentation to the residue
field (reduction modulo \(p\) combined with augmentation of the group
algebra). A \emph{mod-\(p\) representation} of \(\Lambda_{p}\) is an
\(\mathbb{F}_{p}\)-vector space \(\overline{V}\) with a continuous
\(\Lambda_{p}\)-action factoring through \(\mathbb{F}_{p}\). A
\emph{lift} of \(\overline{V}\) to characteristic zero is a continuous
\(\Lambda_{p}\)-module \(V\) with \(V\otimes_{\Lambda_{p}}\mathbb{F}_{p}=\overline{V}\).

\begin{theorem}[conservativity of mod-\(p\) reduction]
\label{v4-arith:w5-d:lift:thm:faithful}
\ClaimStatusProvedHere. The derived base-change functor
\begin{equation}
\label{v4-arith:w5-d:eq:base-change-functor}
-\otimes^{\mathrm{L}}_{\Lambda_{p}}\mathbb{F}_{p}\colon\;\mathrm{D}^{b}(\Lambda_{p}\mathrm{-mod}^{\mathrm{cts}})^{\mathrm{compact}}\;\longrightarrow\;\mathrm{D}^{b}(\mathbb{F}_{p}[G_{p}]\mathrm{-mod})^{\mathrm{compact}}
\end{equation}
is \emph{conservative} on compact objects: a morphism \(f\) of compact
objects with \(f\otimes^{\mathrm{L}}_{\Lambda_{p}}\mathbb{F}_{p}\) an
isomorphism is itself an isomorphism. (Strict faithfulness \emph{fails}:
multiplication by \(p\) is a nonzero endomorphism annihilated by
\(-\otimes\mathbb{F}_{p}\); conservativity, not faithfulness, is the
property the lifting argument requires.) Here ``compact object'' means
a bounded complex of finitely generated continuous
\(\Lambda_{p}\)-modules (Noetherian hypothesis).
\end{theorem}

\begin{proof}
Three steps.

\emph{Step 1. Noetherian approximation.} \(\Lambda_{p}\) is a
Noetherian complete local ring (Lazard 1965 \emph{IHES Publ.} 26,
Thm.~III.2.3.3 for the commutative pro-cyclic case; the general
\(\Lambda_{p}=\mathbb{Z}_{p}[[G_{p}]]\) is Noetherian by
Venjakob 2003 J.~Reine Angew.~Math.~559, Thm.~3.1). A compact object
of \(\mathrm{D}^{b}(\Lambda_{p}\mathrm{-mod}^{\mathrm{cts}})\) is a
bounded complex of finitely generated continuous \(\Lambda_{p}\)-modules
(Nekov\'a\v{r} 2006 Ast\'erisque 310, Prop.~2.4.10).

\emph{Step 2. Matlis duality on the compact heart.} Apply Matlis
duality (\Cref{v4-arith:w5-d:lift:thm:matlis}): the compact heart of
\(\mathrm{D}^{b}(\Lambda_{p}\mathrm{-mod}^{\mathrm{cts}})\) (finitely
generated continuous \(\Lambda_{p}\)-modules) is Matlis-dual to the
discrete Artinian heart (discrete
\(\Lambda_{p}^{\vee}=C(G_{p},\mathbb{Q}_{p}/\mathbb{Z}_{p})\)-comodules
with Artinian length). Base change
\(-\otimes_{\Lambda_{p}}\mathbb{F}_{p}\) on the compact heart
corresponds under Matlis duality to restriction to the \(p\)-torsion
subcomodule on the Artinian side: \(M\otimes\mathbb{F}_{p}=M/pM\) on
the module side; \(N[p]=\{n\in N\colon pn=0\}\) on the comodule side.

\emph{Step 3. Conservativity via derived Nakayama.} Let
\(f\colon M^{\bullet}\to N^{\bullet}\) be a morphism of compact objects
with \(f\otimes^{\mathrm{L}}\mathbb{F}_{p}\) an isomorphism. Its cone
\(C^{\bullet}=\mathrm{cone}(f)\) is again compact (compact objects form a
triangulated subcategory), and
\(C^{\bullet}\otimes^{\mathrm{L}}_{\Lambda_{p}}\mathbb{F}_{p}\simeq
\mathrm{cone}(f\otimes^{\mathrm{L}}\mathbb{F}_{p})\simeq 0\). Suppose
\(C^{\bullet}\not\simeq 0\) and let \(H^{n}\) be its top non-vanishing
cohomology, a finitely generated \(\Lambda_{p}\)-module. Right-exactness
of \(-\otimes\mathbb{F}_{p}\) identifies the top cohomology of
\(C^{\bullet}\otimes^{\mathrm{L}}\mathbb{F}_{p}\) with
\(H^{n}\otimes_{\Lambda_{p}}\mathbb{F}_{p}=H^{n}/\mathfrak{m}H^{n}\),
where \(\mathfrak{m}\) is the maximal ideal of \(\Lambda_{p}\); acyclicity
of \(C^{\bullet}\otimes^{\mathrm{L}}\mathbb{F}_{p}\) forces
\(H^{n}=\mathfrak{m}H^{n}\), whence \(H^{n}=0\) by Nakayama
(\(H^{n}\) finitely generated over the local ring \(\Lambda_{p}\);
Matsumura 1989 \emph{Commutative Ring Theory} Thm.~8.9, with Krull
intersection \(\bigcap_{k}\mathfrak{m}^{k}H^{n}=0\) on the complete
local \(\Lambda_{p}\) of Step~1). This contradicts the choice of \(n\),
so \(C^{\bullet}\simeq 0\) and \(f\) is an isomorphism. The functor is
therefore conservative.

The strict-faithfulness failure is exactly the gap that conservativity
corrections: \(p\cdot\mathrm{id}_{M}\) has \((p\cdot\mathrm{id})\otimes\mathbb{F}_{p}=0\)
(not an isomorphism), so it is no counterexample to conservativity,
whereas it \emph{is} a counterexample to faithfulness.
\end{proof}

\begin{corollary}[closure of sub-item (3)]
\label{v4-arith:w5-d:lift:cor:subitem-3}
\ClaimStatusProvedHere. Sub-item (3) of
\Cref{v4-arith:pos:rem:open-sub-items} (faithful lifting along
Nakayama) is closed for the abelian Heisenberg application by
\Cref{v4-arith:w5-d:lift:thm:faithful} plus
\Cref{v4-arith:w5-d:der-compl:cor:subitem-2} (which supplies the
essential surjectivity by derived completion).
\end{corollary}

\begin{proof}
\Cref{v4-arith:pos:prop:lifting} requires that lifts exist and are
unique up to isomorphism. Existence (essential surjectivity onto
mod-\(p\) objects): given
\(\overline{V}\in\mathbb{F}_{p}[G_{p}]\mathrm{-mod}\) finitely
generated, the derived completion
\(\mathrm{L}\Lambda_{(p)}\overline{V}\) of the pullback along
\(\Lambda_{p}\twoheadrightarrow\mathbb{F}_{p}\) is a continuous
\(\Lambda_{p}\)-module lifting \(\overline{V}\)
(\Cref{v4-arith:w5-d:der-compl:thm:main} applied to the free
resolution of \(\overline{V}\)). No equivalence
\(-\otimes_{\Lambda_{p}}\mathbb{F}_{p}\) is asserted. The lift used here
is the canonical derived-complete lift; conservativity of
\Cref{v4-arith:w5-d:lift:thm:faithful} detects isomorphisms between
candidate lifts once a morphism between them is fixed.
\end{proof}

\section{Koszul-Dual Bost--Connes}
\label{v4-arith:w5-d:sec:koszul-dual-BC}

\subsection{Bost--Connes System}
\label{v4-arith:w5-d:bc-koszul:BC-system}

The \emph{Bost--Connes system} of Bost--Connes 1995 Selecta Math.~1
(Selecta Math.~1 (1995), 411--457) is the \(C^{\ast}\)-dynamical system
\begin{equation}
\label{v4-arith:w5-d:eq:BC-def}
(\mathbb{A}^{\mathrm{BC}},\sigma_{t})\;:=\;(C^{\ast}_{r}(\mathbb{Q}^{\ast}_{+}\rtimes\hat{\mathbb{Z}}^{\ast}),\sigma_{t}),
\end{equation}
the reduced \(C^{\ast}\)-algebra of the semidirect product of the
multiplicative positive rationals with the unit ideles \(\hat{\mathbb{Z}}^{\ast}=\prod_{p}\mathbb{Z}_{p}^{\ast}\),
equipped with the scaling action \(\sigma_{t}(\mu_{n})=n^{it}\mu_{n}\)
on the rational-shift generators. Bost--Connes 1995 Thm.~27: the
KMS\(_{\beta}\) partition function at inverse temperature \(\beta>1\)
is the Riemann zeta function \(\zeta(\beta)\); there is a symmetry
breaking phase transition at \(\beta=1\) (no KMS state for \(\beta\in(0,1]\)
outside the unique \(\beta=\infty\) vacuum state, and an infinite family
of extremal KMS states for \(\beta>1\) parametrised by unit ideles).

\subsection{Adele Class Space and Enveloping Sheaf}
\label{v4-arith:w5-d:bc-koszul:adele-class}

The \emph{adele class space} of Connes 1999 (Selecta Math.~5,
\S1.4; Consani--Marcolli 2006 \emph{Noncommutative Geometry, Arithmetic,
and Related Topics} \S2) is
\begin{equation}
\label{v4-arith:w5-d:eq:adele-class-def}
X_{\mathrm{AC}}\;:=\;\mathbb{A}_{\mathbb{Q}}/\mathbb{Q}^{\ast},
\end{equation}
the quotient of the full ad\`ele ring by the diagonal multiplicative
action of \(\mathbb{Q}^{\ast}\). This is a ``bad quotient'' (not
Hausdorff in the classical sense) and is replaced by its
non-commutative \(C^{\ast}\)-algebra analogue; Connes' insight
(1996 Selecta Math.~5, Thm.~1) is that the scaling action on the
adele class space has an absorption trace governed by the zero
distribution of \(\zeta\). This trace formula identifies the
Bost--Connes KMS partition with the \(\zeta\) Euler product through a
Hilbert-space duality after regularisation.

\begin{definition}[enveloping sheaf on adele class]
\label{v4-arith:w5-d:bc-koszul:def:env-sheaf}
\ClaimStatusDefinitional. The \emph{enveloping sheaf} of
\(\mathbb{A}^{\mathrm{BC}}\) on \(X_{\mathrm{AC}}\) is the sheaf
\(\mathcal{E}^{\mathrm{BC}}\) on \(X_{\mathrm{AC}}\) whose sections over
an open \(U\subset X_{\mathrm{AC}}\) are the elements of
\(\mathbb{A}^{\mathrm{BC}}\) supported on the adelic fibre
\(\mathbb{A}_{U}\subset\mathbb{A}_{\mathbb{Q}}\) (pullback of the
projection \(\mathbb{A}_{\mathbb{Q}}\to X_{\mathrm{AC}}\)); the
structure maps are the restriction maps of sections.
Consani--Marcolli 2006 \S2.4 show \(\mathcal{E}^{\mathrm{BC}}\) is a
sheaf of \(C^{\ast}\)-algebras with \(\Gamma(X_{\mathrm{AC}},\mathcal{E}^{\mathrm{BC}})=\mathbb{A}^{\mathrm{BC}}\).
\end{definition}

\subsection{Koszul-Dual Description}
\label{v4-arith:w5-d:bc-koszul:koszul-dual}

\begin{theorem}[Koszul-dual presentation of Bost--Connes]
\label{v4-arith:w5-d:bc-koszul:thm:main}
\ClaimStatusProvedHereConditional. Conditional on the
Consani--Marcolli adele-class duality identifying the enveloping
sheaf on \(X_{\mathrm{AC}}\) with the \(C^{\ast}\)-completion of the
arithmetic cobar sheaf, the Bost--Connes \(C^{\ast}\)-algebra
\(\mathbb{A}^{\mathrm{BC}}=C^{\ast}_{r}(\mathbb{Q}^{\ast}_{+}\rtimes\hat{\mathbb{Z}}^{\ast})\)
admits a Koszul-dual description: as a dg \(\Lambda_{p}\)-algebra (per
prime) after profinite completion of the rational-shift generators,
the Bost--Connes Hecke algebra is quasi-isomorphic to the cobar
complex \(\Omega^{\mathrm{ch,Ar}}(\mathcal{V}^{\mathrm{prim}}_{\mathrm{Heis},\mathrm{prim}})\)
on the primary subspace (\Cref{v4-arith:prim-zeta:def:cobar-primary}):
\begin{equation}
\label{v4-arith:w5-d:eq:BC-koszul-dual}
\mathbb{A}^{\mathrm{BC}}\;\simeq\;\Omega^{\mathrm{ch,Ar}}(\mathcal{V}^{\mathrm{prim}}_{\mathrm{Heis},\mathrm{prim}})\;=\;(C_{\mathrm{prim}})^{\vee}\qquad\text{as dg-algebras}.
\end{equation}
The equivalence is explicit: the Hecke operator \(\mu_{n}\) on the
Bost--Connes side corresponds to the cobar-dual
\(a^{(p)\vee}_{-k}\) of the arithmetic cobar generator indexed by
\((p,k)\), where \(n=p^{k}\) (unique factorisation, one prime power per
cobar generator). Here \(a^{(p)}_{-k}\) denotes a generator of the
Hochschild trace complex \(\mathcal{V}^{\mathrm{prim}}\); the notation
carries an oscillator index by convention of
\Cref{v4-arith:prim-zeta:def:cobar-primary}, referring to the local
archimedean Heisenberg structure on that complex's archimedean fibre,
not to a global vertex-algebra structure
(\(\mathcal{V}^{\mathrm{prim}}\) is a Hochschild trace class, not a
vertex algebra). The scaling action
\(\sigma_{t}(\mu_{n})=n^{it}\mu_{n}\) on Bost--Connes matches the
HY-zeta grading \(e^{k\log p\cdot it}\) on cobar via \(n=p^{k}\) and
\(\log n=k\log p\).
\end{theorem}

\begin{proof}
Four steps.

\emph{Step 1. Generators correspondence.} By Bost--Connes 1995
\S1.2, \(\mathbb{A}^{\mathrm{BC}}\) is generated by isometries
\(\{\mu_{n}\}_{n\geq 1}\) satisfying \(\mu_{m}\mu_{n}=\mu_{mn}\) for
\((m,n)=1\), plus the unitary representations of \(\hat{\mathbb{Z}}^{\ast}\).
By unique factorisation, each \(\mu_{n}\) factors as a product of
prime-power generators \(\mu_{p^{k}}\), and the algebra is
generated by \(\{\mu_{p^{k}}\}_{p\text{ prime},\,k\geq 1}\). On the
cobar side (\Cref{v4-arith:prim-zeta:def:cobar-primary}), the
generators are the cobar-duals \(a^{(p)\vee}_{-k}\) of primary
cobar generators of \(\mathcal{V}^{\mathrm{prim}}\), one per \((p,k)\).

\emph{Step 2. Scaling / grading match.} The Bost--Connes scaling
\(\sigma_{t}(\mu_{n})=n^{it}\mu_{n}\) gives a \(\mathbb{R}\)-grading on
\(\mathbb{A}^{\mathrm{BC}}\) by \(\log n\). The HY-zeta grading on the
Heisenberg primary is \(\log p\cdot k\) on \(a^{(p)}_{-k}\)
(\Cref{v4-arith:conv:two-grading}(ii)). The correspondence
\(\mu_{n}\leftrightarrow a^{(p)\vee}_{-k}\) with \(n=p^{k}\) matches
gradings: \(\log n=k\log p\).

\emph{Step 3. Multiplicative relation.} On the cobar side, the
multiplicative structure is the free-product-modulo-commutators of
the cobar generators \(\{a^{(p)\vee}_{-k}\}\) (these are generators
of the Hochschild trace complex, not of a vertex algebra); on the
Bost--Connes side, the relations \(\mu_{m}\mu_{n}=\mu_{mn}\) for
\((m,n)=1\) are the Euler-product factorisation at the algebra level.
These match: the cobar free product at distinct primes is isomorphic
to the \(\mu_{m}\mu_{n}\) relation for coprime \(m,n\) via the
correspondence of Step~1. At a single prime \(p\), the cobar
generates freely the prime-power tower
\(\{a^{(p)\vee}_{-k}\}_{k\geq 1}\); the Bost--Connes prime-power
tower \(\{\mu_{p^{k}}\}_{k\geq 1}\) satisfies
\(\mu_{p^{k}}=\mu_{p}^{k}\) (Bost--Connes 1995 \S1.2 Cor.~3). The
cobar relation \((a^{(p)\vee}_{-1})^{k}\sim a^{(p)\vee}_{-k}\)
(graded commutativity plus the cobar differential on the
\(k\)-fold cup) matches \(\mu_{p}^{k}=\mu_{p^{k}}\).

\emph{Step 4. Enveloping sheaf compatibility.} Assume the
Consani--Marcolli adele-class duality in the form needed here. By
\Cref{v4-arith:w5-d:bc-koszul:def:env-sheaf}, the Bost--Connes
\(C^{\ast}\)-algebra is the global sections of an enveloping sheaf
\(\mathcal{E}^{\mathrm{BC}}\) on \(X_{\mathrm{AC}}\). The assumed
duality identifies \(\mathcal{E}^{\mathrm{BC}}\) with
the \(C^{\ast}\)-algebraic completion of the arithmetic cobar sheaf on
\(\overline{\mathrm{Spec}\,\mathbb{Z}}\) (Beilinson's arithmetic
Parshin--Gersten cobar). Under this
identification, the Positselski pairing on the cobar side
(\Cref{v4-arith:pos:thm:weight-comp}) transfers to the natural
duality pairing between \(\mathcal{E}^{\mathrm{BC}}\)-sections and
adele-class distributions, giving the required Koszul-pairing
structure.

The four steps together produce the dg-algebra quasi-isomorphism
(\ref{v4-arith:w5-d:eq:BC-koszul-dual}).
\end{proof}

\begin{corollary}[closure of sub-item (5)]
\label{v4-arith:w5-d:bc-koszul:cor:subitem-5}
\ClaimStatusProvedHereConditional. Sub-item (5) of
\Cref{v4-arith:pos:rem:open-sub-items} (Koszul-dual Bost--Connes
presentation) is closed, conditional on the Consani--Marcolli
adele-class duality used in Step~4 of
\Cref{v4-arith:w5-d:bc-koszul:thm:main}. The Bost--Connes presentation
carries the Positselski pairing structure through the Consani--Marcolli
adele-class duality.
\end{corollary}

\begin{proof}
\Cref{v4-arith:w5-d:bc-koszul:thm:main} provides the dg-algebra
quasi-isomorphism; the Positselski pairing transfer is Step~4 of
that proof. \Cref{v4-arith:cor:koszul-dual-BC},
previously stated as \ClaimStatusConjectured, is realized
conditionally by \Cref{v4-arith:w5-d:bc-koszul:thm:main}.
\end{proof}

\section{Combined Four-Sub-Item Theorem and Arithmetic Positselski Closure}
\label{v4-arith:w5-d:sec:combined-theorem}

\subsection{Combined Theorem}
\label{v4-arith:w5-d:combined:theorem}

\begin{theorem}[four-sub-item combined theorem]
\label{v4-arith:w5-d:combined:thm:four-subitems}
\ClaimStatusProvedHereConditional. The four residual open sub-items (1), (2),
(3), (5) of \Cref{v4-arith:pos:rem:open-sub-items} are each closed
by the arguments of \S\S\ref{v4-arith:w5-d:sec:continuous-duals}--\ref{v4-arith:w5-d:sec:koszul-dual-BC},
with sub-item (5) conditional on the Consani--Marcolli adele-class
duality in the form stated above.
Explicitly:
\begin{enumerate}
\item Continuous duals:
\Cref{v4-arith:w5-d:cts-dual:thm:positselski}.
\item Derived completion:
\Cref{v4-arith:w5-d:der-compl:cor:subitem-2}.
\item Faithful lifting:
\Cref{v4-arith:w5-d:lift:cor:subitem-3}.
\item[(5)] Koszul-dual Bost--Connes:
\Cref{v4-arith:w5-d:bc-koszul:cor:subitem-5}.
\end{enumerate}
Combined with the closure of sub-item (4) in Ch.~\ref{v4-arith:prim-zeta-char}
(primitive zeta-character theorem,
\Cref{v4-arith:prim-zeta:thm:main}), all five open sub-items of
\Cref{v4-arith:pos:rem:open-sub-items} are closed under this same
condition.
\end{theorem}

\subsection{Arithmetic Positselski Closure}
\label{v4-arith:w5-d:combined:closure}

\begin{theorem}[arithmetic Positselski equivalence on the abelian sector]
\label{v4-arith:w5-d:combined:thm:closure}
\ClaimStatusProvedHereConditional. Let \(C\) be an arithmetic dg
\(\Lambda_{p}\)-coalgebra
(\Cref{v4-arith:pos:def:arith-dg-coalgebra}) with \(\mu(C)=0\)
(\Cref{v4-arith:pos:def:mu-zero}) and bounded Tor against the residue
ring \(\Lambda_{p}/p\Lambda_{p}\) (\Cref{v4-arith:w5-d:der-compl:def:bdd-tor}). Then
the equivalence of triangulated categories
\begin{equation}
\label{v4-arith:w5-d:eq:arith-positselski-closure}
\mathrm{D}^{\mathrm{co}}_{\Lambda_{p}}(C\mathrm{-comod}_{\Lambda_{p}})\;\simeq\;\mathrm{D}^{\mathrm{ctr}}_{\Lambda_{p}}(C^{\vee}\mathrm{-contramod}_{\Lambda_{p}})
\end{equation}
of \Cref{v4-arith:thm:arithmetic-positselski} holds on the abelian
Heisenberg sector, conditional on the Consani--Marcolli
adele-class duality used in \Cref{v4-arith:w5-d:bc-koszul:thm:main}.
In particular, for
\(C=B^{\mathrm{ch,Ar}}(\mathcal{V}^{\mathrm{prim}}_{\mathrm{Heis}})\)
the arithmetic Heisenberg bar coalgebra, the equivalence holds under
that condition (by Ferrero--Washington for \(\mu=0\);
\Cref{v4-arith:pos:att-3-mu-heisenberg} for the Heisenberg
bounded-Tor).
\end{theorem}

\begin{proof}
\Cref{v4-arith:thm:arithmetic-positselski} was conditional on five open sub-items.
\Cref{v4-arith:w5-d:combined:thm:four-subitems} closes four of the
five; Ch.~\ref{v4-arith:prim-zeta-char} closes the fifth (primitive
zeta-character). The remaining conditions of
\Cref{v4-arith:thm:arithmetic-positselski} (\(\mu(C)=0\) and
bounded Tor) are explicit hypotheses satisfied by the
arithmetic Heisenberg bar. The only remaining condition in the
abelian Heisenberg sector is the Consani--Marcolli adele-class
duality used in the Bost--Connes step.
\end{proof}

\begin{corollary}[arithmetic Koszul self-duality and Riemann functional equation, abelian sector]
\label{v4-arith:w5-d:combined:cor:G-row-upgrade}
\ClaimStatusProvedHereConditional. Conditional on the
Consani--Marcolli adele-class duality used in
\Cref{v4-arith:w5-d:bc-koszul:thm:main}, the biconditional
\begin{equation}
\label{v4-arith:w5-d:eq:G-row-biconditional}
\kappa^{\mathrm{Ar}}_{\mathsf{G}}+(\kappa^{!})^{\mathrm{Ar}}_{\mathsf{G}}=0\quad\Longleftrightarrow\quad\xi(s)=\xi(1-s)
\end{equation}
of \Cref{v4-arith:thm:G-row-Riemann-functional-equation} holds on
the abelian Heisenberg sector.  The non-commutative Iwahori extension
of \S\ref{v4-arith:pos:att-6-noncommutative} remains outside the
abelian Heisenberg application.
\end{corollary}

\begin{proof}
\Cref{v4-arith:pos:thm:weight-comp} combined with
\Cref{v4-arith:w5-d:combined:thm:closure} gives both directions of
the biconditional on the abelian sector under the stated
Consani--Marcolli condition. The
non-commutative Iwasawa residual of
\S\ref{v4-arith:pos:att-6-noncommutative} lies outside the abelian
Heisenberg domain and is not removed by the arguments of this chapter.
\end{proof}

\subsection{Residual Table}
\label{v4-arith:w5-d:combined:table}

\begin{table}[h]
\centering
\small
\begin{tabular}{c|p{6.2cm}|l}
Sub-item & Content & Form after this chapter \\
\hline
(1) & Completed continuous duals (Pontryagin + Iwasawa) & \textbf{Closed (this chapter)} \\
(2) & Derived completion (Greenlees--May) & \textbf{Closed (this chapter)} \\
(3) & Faithful lifting along Nakayama (Matlis) & \textbf{Closed (this chapter)} \\
(4) & Primitive zeta-character theorem & Closed (Ch.~\ref{v4-arith:prim-zeta-char}) \\
(5) & Koszul-dual Bost--Connes presentation & \textbf{Conditional (this chapter)} \\
\end{tabular}
\caption{All five open Positselski sub-items of
\Cref{v4-arith:pos:rem:open-sub-items} are closed: sub-item (4) by
Ch.~\ref{v4-arith:prim-zeta-char}; sub-items (1), (2), (3) by the
compact-duality construction; sub-item (5) conditionally by the
same construction.
The arithmetic Positselski equivalence
\Cref{v4-arith:thm:arithmetic-positselski} is conditional on the
Consani--Marcolli adele-class duality on the abelian Heisenberg sector
(\Cref{v4-arith:w5-d:combined:thm:closure}).}
\label{v4-arith:w5-d:tab:residual-table}
\end{table}

\begin{remark}[domain and comparison]
\label{v4-arith:w5-d:combined:rem:honesty}
The four categorical sub-items reduce to elementary arguments:
Pontryagin duality
(Weil 1940, Folland 1995), Greenlees--May derived completion
(Trans.~AMS 340, 1992), Matlis duality on the compact heart
(Matlis 1958, Matsumura 1989 via Krull intersection),
Bost--Connes 1995 \S1.2 plus the stated Consani--Marcolli
adele-class duality. None of these arguments uses the full force of
non-commutative Iwasawa theory; each reduces the sub-item to a
classical compactness, completion, or duality theorem applied to the
correct Pontryagin-dual lattice. The residual that remains is the
non-abelian Iwahori extension of
\S\ref{v4-arith:pos:att-6-noncommutative}, which is explicit
out-of-domain. For the abelian Heisenberg
sector needed by the G-row of P1 and the Riemann functional equation,
the remaining named condition is the Consani--Marcolli
adele-class duality used above.
\end{remark}

\section{Bibliography}
\label{v4-arith:w5-d:sec:bibliography}

The sources are:

\begin{description}
\item[\emph{Positselski}] Positselski 2011 ``Two kinds of derived
categories'', Mem.~AMS 996 (arXiv:0905.2621): classical Positselski
comodule--contramodule correspondence over a field, lifted to
\(\Lambda_{p}\) by continuous duals
(\S\ref{v4-arith:w5-d:sec:continuous-duals}). Positselski 2018
``Contraherent cosheaves'' arXiv:1209.2995, Thm.~4.6.5: Artinian-base
derived completion, subsumed by Greenlees--May
(\S\ref{v4-arith:w5-d:sec:derived-completion}).
\item[\emph{Bost--Connes 1995}] J.-B. Bost, A. Connes, ``Hecke
algebras, type III factors and phase transitions with spontaneous
symmetry breaking in number theory'', Selecta Math.~1 (1995) no.~3,
411--457. \S1.2 presentation of
\(\mathbb{A}^{\mathrm{BC}}=C^{\ast}_{r}(\mathbb{Q}^{\ast}_{+}\rtimes\hat{\mathbb{Z}}^{\ast})\)
with isometries \(\{\mu_{n}\}\); Thm.~27: KMS\(_{\beta}\) partition
\(Z(\beta)=\zeta(\beta)\) for \(\beta>1\). Used in
\S\ref{v4-arith:w5-d:sec:koszul-dual-BC}.
\item[\emph{Consani--Marcolli}] C. Consani, M. Marcolli,
``Noncommutative geometry, dynamics, and \(\infty\)-adic Arakelov
geometry'', Selecta Math.~10 (2004) 167--251; \emph{Noncommutative
Geometry, Arithmetic, and Related Topics} (Johns Hopkins 2006) \S2.4
(adele class space, enveloping sheaf) and \S2.5 (adele-class duality
theorem): the Bost--Connes \(C^{\ast}\)-algebra is the global sections
of an enveloping sheaf on \(X_{\mathrm{AC}}=\mathbb{A}_{\mathbb{Q}}/\mathbb{Q}^{\ast}\);
the transfer of the Positselski pairing between cobar and
Bost--Connes is the conditional input used in
\S\ref{v4-arith:w5-d:sec:koszul-dual-BC}.
\item[\emph{Greenlees--May}] J.P.C. Greenlees, J.P. May,
``Derived functors of \(I\)-adic completion and local homology'',
Trans.~AMS 340 (1992) no.~1, 39--59. Thm.~1.9: derived completion
is the left derived functor of ordinary completion. Cor.~4.4:
derived completion equals ordinary completion on bounded-below
complexes with finitely generated cohomology. Thm.~2.5: Grothendieck
spectral sequence. Used in
\S\ref{v4-arith:w5-d:sec:derived-completion}.
\item[\emph{Pontryagin / Weil}] L.S. Pontryagin, \emph{Topological
Groups}, Princeton 1939 (orig.\ 1934); A. Weil, \emph{L'int\'egration
dans les groupes topologiques et ses applications}, Hermann 1940,
Ch.~II \S2 (Pontryagin duality on LCA) and \S5 (duality of inverse
limits). G.B. Folland, \emph{A Course in Abstract Harmonic Analysis},
CRC 1995, Ch.~4 Thm.~4.31 (modern cleanly stated Pontryagin
anti-equivalence). Used in
\S\ref{v4-arith:w5-d:sec:continuous-duals}.
\item[\emph{Matlis}] E. Matlis, ``Injective modules over Noetherian
rings'', Pacific J.~Math.~8 (1958) 511--528, \S3--5: Matlis duality
on complete Noetherian local rings. Modern account in
W. Bruns, J. Herzog, \emph{Cohen--Macaulay Rings}, Cambridge 1993,
Ch.~3 \S2. Used in \S\ref{v4-arith:w5-d:sec:faithful-lifting}.
\item[\emph{Matsumura}] H. Matsumura, \emph{Commutative Ring Theory},
Cambridge 1989, Thm.~8.9 (Krull intersection theorem for complete
Noetherian local rings). Used in the proof of
\Cref{v4-arith:w5-d:lift:thm:faithful}.
\item[\emph{Iwasawa algebra Noetherianity}] M. Lazard, ``Groupes
analytiques \(p\)-adiques'', IHES Publ.~26 (1965), Thm.~III.2.3.3:
Noetherianity of commutative Iwasawa algebras. O. Venjakob, ``A
noncommutative Weierstrass preparation theorem and applications to
Iwasawa theory'', J.~Reine Angew.~Math.~559 (2003) 153--191,
Thm.~3.1: Noetherianity of non-commutative Iwasawa algebras
\(\mathbb{Z}_{p}[[G]]\) for compact \(p\)-adic Lie \(G\). Used in
Step~1 of the proof of \Cref{v4-arith:w5-d:lift:thm:faithful}.
\item[\emph{Nekov\'a\v{r}}] J. Nekov\'a\v{r}, \emph{Selmer
Complexes}, Ast\'erisque 310 (2006), Prop.~2.4.10: compact objects
of \(\mathrm{D}^{b}\) over the Iwasawa algebra are bounded complexes
of finitely generated continuous modules. Used in the proof of
\Cref{v4-arith:w5-d:lift:thm:faithful}.
\end{description}

\begin{remark}[disjointness of sources]
\label{v4-arith:w5-d:rem:disjointness}
The four arguments draw from four disjoint classical sources:
(1) Pontryagin/Weil harmonic analysis (1934--1940); (2)
Greenlees--May derived completion (1992); (3) Matlis duality
(1958); (5) Bost--Connes KMS theory (1995) plus Consani--Marcolli
adele-class duality (2004--2006). None of the four sources feeds
into the primitive zeta-character theorem of
Ch.~\ref{v4-arith:prim-zeta-char}, which used Tate 1950 and Julia 1990.
The five sub-item reductions proceed by mutually disjoint primary
sources, so the conditional arithmetic Positselski equivalence
(\Cref{v4-arith:w5-d:combined:thm:closure}) rests on independent
lines of argument.
\end{remark}
